% TOP_PROBLEM: {"id": 30007057, "title": "Numerical semiampleness of projective adjoints", "category_id": 5, "category": {"id": 5, "name": "algebraic_geometry", "display_name": "Algebraic Geometry", "description": "Geometric objects defined by polynomial equations.", "slug": "algebraic-geometry", "order_index": 5, "created_at": "2026-07-31T15:26:25.671Z"}, "status": "open"}
% FIRST_POSTED: 2026-10-06
% ATTEMPT_STATUS: solved
% OPENAI_CLAIM: {"schema_version":1,"source_repo":"https://github.com/openai/math","source_commit":"adc7f1241b42e322a6451854ab7e4b4c146bf78a","release_date":"2026-10-06","match":"full","resolution":"proved","independently_reviewed":false,"second_pass":"agreed","families":[{"family":"036","title":"Numerical semiampleness and generalized minimal models","subject":"Algebraic and complex geometry","manuscripts":[{"title":"Numerical Semiampleness of Nef Adjoint Divisors","pdf_path":"preprints/Numerical-Semiampleness-of-Nef-Adjoint-Divisors-October-3-2026/paper.pdf","readme_path":"preprints/Numerical-Semiampleness-of-Nef-Adjoint-Divisors-October-3-2026/README.md","date":"2026-10-03","theorem_ref":"preprints/Numerical-Semiampleness-of-Nef-Adjoint-Divisors-October-3-2026/build/sections/introduction.tex#thm:main"}],"lean":null,"reasoning_trace":null}]}
\documentclass[11pt]{article}
\usepackage[margin=25mm]{geometry}
\usepackage{fontspec}
\setmainfont{Latin Modern Roman}
\usepackage{amsmath,amssymb,mathtools}
\usepackage{unicode-math}
\setmathfont{Latin Modern Math}
\usepackage{enumitem}
\usepackage{xurl,hyperref,bookmark}
\hypersetup{colorlinks=true,urlcolor=blue}
\setcounter{secnumdepth}{0}
\setlength{\parindent}{0pt}
\setlength{\parskip}{5pt}
\setlength{\emergencystretch}{3em}
\begin{document}
\section*{Numerical semiampleness of projective adjoints}

\subsection{OpenAI's claim}
For projective characteristic-zero klt pairs, a nef ordinary adjoint plus a nef rational divisor is claimed to be numerically equivalent to a semiample divisor whenever the ordinary adjoint is pseudo-effective. The cited manuscript argues for that result; its relation to this problem is specified below.

\subsection{Outline of the argument}
Manuscript 1 (family 036, 2026-10-03): Ordinary log abundance and terminating programs supply the birational reductions. The remaining full-nef-dimension Calabi--Yau case is treated with jet and Frobenius estimates, after which the positive part descends as a semiample numerical representative; flat nontorsion bundles explain why numerical equivalence is the conclusion.

\subsection{Scope relative to this problem}
Manuscript 1 (family 036, 2026-10-03): The problem selects the mathematical statement summarized here, including its stated hypotheses and parameter range. The cited theorem gives a claimed proof of that assertion in full.

\subsection{Formal verification}
Family 036 has no Lean scope page in this snapshot; no formalization of these manuscript results is listed.

This site has not executed the Lean files or checked their dependencies. These descriptions report the scope of the linked declarations, rather than a new verification of the proofs.

\subsection{Status caveat}
OpenAI attributes these results to an internal model. Its release notice cautions that results lacking full formalization can contain mistakes. This site has not independently assessed the proofs; a solved label records the scope of the reported claim.

\subsection{Sources}
\begin{itemize}
\item OpenAI, \emph{Numerical Semiampleness of Nef Adjoint Divisors}, 2026-10-03; citation key \texttt{OAI:Numerical-Semiampleness-of-Nef-Adjoint-Divisors-October-3-2026}. \href{https://github.com/openai/math/blob/adc7f1241b42e322a6451854ab7e4b4c146bf78a/preprints/Numerical-Semiampleness-of-Nef-Adjoint-Divisors-October-3-2026/paper.pdf}{PDF (GitHub)}; \href{https://github.com/openai/math/raw/adc7f1241b42e322a6451854ab7e4b4c146bf78a/preprints/Numerical-Semiampleness-of-Nef-Adjoint-Divisors-October-3-2026/paper.pdf}{PDF (direct)}; \href{https://github.com/openai/math/blob/adc7f1241b42e322a6451854ab7e4b4c146bf78a/preprints/Numerical-Semiampleness-of-Nef-Adjoint-Divisors-October-3-2026/README.md}{Manuscript page and citation}; \href{https://github.com/openai/math/blob/main/preprints/Numerical-Semiampleness-of-Nef-Adjoint-Divisors-October-3-2026/paper.pdf}{Latest version}.
\item \href{https://github.com/openai/math/blob/adc7f1241b42e322a6451854ab7e4b4c146bf78a/CONTENTS.md}{OpenAI catalogue entry}.
\end{itemize}
\end{document}
